% -------------------------------------------------
% style.tex – Ortak paketler ve tcolorbox tanımları
% -------------------------------------------------
\usepackage[utf8]{inputenc}
\usepackage[turkish]{babel}
\usepackage{amsmath,amssymb,amsfonts,amsthm}
\usepackage{geometry}
\geometry{top=2.5cm,bottom=2.5cm,left=2.5cm,right=2.5cm}
\usepackage{xcolor}
\usepackage{tikz,pgfplots}
\pgfplotsset{compat=1.18}
\usepackage[most]{tcolorbox}
\usepackage{hyperref}
\usepackage{caption,booktabs,tabularx}

% ---------- Renk tanımları ----------
\definecolor{purp}{RGB}{102,51,153}
\definecolor{tealblue}{RGB}{0,120,130}
\definecolor{mainblue}{RGB}{26,82,118}
\definecolor{alertred}{RGB}{180,40,40}
\definecolor{okgreen}{RGB}{34,139,34}

% ---------- Tek‑argümanlı kutular ----------
\newtcolorbox{box}[2][]{
  colback=white,
  colframe=purp,
  title={#2},
  arc=2mm,
  breakable,
  #1
}

\newtcolorbox{uyari}[1][]{
  colback=yellow!10,
  colframe=orange,
  title={Uyarı},
  arc=2mm,
  breakable,
  #1
}

% ---------- Diğer yardımcı kutular (isteğe bağlı) ----------
\newtcolorbox{ornek}[2][]{
  colback=mainblue!5!white,
  colframe=mainblue,
  title={#2},
  arc=2mm,
  breakable,
  #1
}

\newtcolorbox{odev}[2][]{
  colback=okgreen!5!white,
  colframe=okgreen,
  title={Çözümü Olmayan Ödev ve Tam Çözümü: #2},
  arc=2mm,
  breakable,
  #1
}

\newtcolorbox{tanim}[2][]{
  colback=gray!5!white,
  colframe=gray!60!black,
  title={#2},
  arc=2mm,
  breakable,
  #1
}

\newtcolorbox{sezgi}[2][]{
  colback=tealblue!5!white,
  colframe=tealblue,
  title={Sezgisel Bakış ve Basit Anlatım: #2},
  arc=2mm,
  breakable,
  #1
}
